\documentclass[11pt]{article}

\usepackage{amsmath,amssymb,amsfonts}
\usepackage[margin=1in]{geometry}
\usepackage{hyperref}
\usepackage{bm}
\usepackage{booktabs}

\hypersetup{colorlinks=true, linkcolor=blue, citecolor=blue, urlcolor=blue}

\newcommand{\ee}{\bm{e}}
\newcommand{\BoW}{\mathrm{BoW}}
\newcommand{\Tr}{\mathrm{Tr}}

\usepackage{xcolor}
\newcommand{\claude}[1]{{\color{red}\textsf{#1}}}
\newcommand{\old}[1]{{\color{gray}#1}}
\newcommand{\lcq}[1]{\textcolor{blue}{\textbf{[LCQ: #1]}}}

\title{Sequential Dependence from Nearest-Neighbor Perturbation:\\Dressed Embeddings as Bigram-Weighted Averages}
\author{}
\date{}

\begin{document}
\maketitle

\begin{abstract}
We perturb the bag-of-words energy model by a nearest-neighbor interaction $\lambda \sum_t \ee_{w_t} \cdot \ee_{w_{t+1}}$ and work to first order in $\lambda$. Absorbing the perturbation into the embeddings defines \emph{dressed} embeddings $\ee_i'$ whose shift from the bare values is a bigram-weighted average of all other embeddings: $\delta \ee_i = (\lambda/2)\sum_j P_{ij}^{(b)}\, \tilde{\ee}_j$, where $P_{ij}^{(b)}$ is the bigram transition probability determined self-consistently by the model.
\end{abstract}

\tableofcontents

\section{The perturbed model}\label{sec:setup}

\subsection{Notation}

Consider a vocabulary of $N$ tokens and a context window of length $L$. A sequence is $(w_1, w_2, \ldots, w_t, \ldots, w_L)$ where $t$ labels the context position. For token type $i \in \{1, \ldots, N\}$, define the count fluctuation
\begin{equation}\label{eq:count}
    n_i = \sum_{t=1}^{L} \delta_{i, w_t} - P_i\,,
\end{equation}
where $P_i$ is the expected frequency of token $i$. The count covariance is $\Sigma_{ij} = \langle n_i\, n_j \rangle$, and the precision (inverse covariance) matrix is $G = \Sigma^{-1}$.

\subsection{Unperturbed theory}

The bag-of-words model assigns
\begin{equation}\label{eq:bow}
    p_{\BoW}(w_1, \ldots, w_L) \propto \exp\!\Big(-\sum_{i,j} G_{ij}\, n_i\, n_j\Big)\,,
\end{equation}
with $G_{ij} = (\Sigma^{-1})_{ij}$. The energy is a quadratic form in the count fluctuations, invariant under permutations of positions. The central result of the unperturbed theory is that the optimal embedding vectors satisfy
\begin{equation}\label{eq:bow-result}
    \ee_i \cdot \ee_j = \Sigma_{ij}\,,
\end{equation}
i.e., embedding dot products equal the count covariance.

\subsection{Nearest-neighbor perturbation}

We add a nearest-neighbor interaction in position space:
\begin{equation}\label{eq:perturbed}
    \boxed{\;p(w_1, \ldots, w_L) \propto \exp\!\Big(-\sum_{i,j} G_{ij}\, n_i\, n_j - \lambda \sum_{t=1}^{L-1} \ee_{w_t} \cdot \ee_{w_{t+1}}\Big)\;}
\end{equation}
\lcq{I think the perturbation should be centered by average of $e$ also?}
The coupling $\lambda$ controls the strength of the sequential dependence. We work in the regime $\lambda \ll 1$ with $\|\ee_j\| = O(1)$, so the perturbation is small. At $\lambda = 0$ we recover the permutation-invariant BoW model.

\section{Bilinear architecture}\label{sec:arch}

\subsection{Embedding and centering}

Expanding the energy~\eqref{eq:perturbed} in position sums, the full log-probability (up to sequence-independent constants) is $\log p = -\sum_{t,t'} G_{w_t, w_{t'}} - \lambda \sum_{t} \ee_{w_t} \cdot \ee_{w_{t+1}} + \text{const}$. This suggests a three-step architecture:
\begin{enumerate}
    \item \textbf{Embed.} Map each token $w_t = i$ to its embedding vector $\ee_i \in \mathbb{R}^d$.
    \item \textbf{Center.} Apply an affine transformation:
    \begin{equation}\label{eq:center}
        \tilde{\bm{x}}_t = A^{-1}\big(\ee_{w_t} - \bar{\ee}\big)\,,
    \end{equation}
    where $\bar{\ee} = \sum_j P_j\, \ee_j$ is the mean embedding and $A$ is a (diagonal) scaling matrix.
    \item \textbf{Bilinear product.} Compute the quadratic form
    \begin{equation}\label{eq:bilinear}
        \log p \;\propto\; \tilde{\bm{x}}^\top W\, \tilde{\bm{x}}\,,
    \end{equation}
    where $\tilde{\bm{x}} = (\tilde{\bm{x}}_1, \ldots, \tilde{\bm{x}}_L)^\top$ is the $L \times d$ matrix of centered embeddings and $W$ is an $L \times L$ matrix acting in context space.
\end{enumerate}

\subsection{Structure of $W$}

The BoW part contributes an all-to-all coupling (every position pair $(t, t')$ gets weight $G_{w_t, w_{t'}}$). The nearest-neighbor part contributes only to adjacent pairs. The combined weight matrix in context space is
\begin{equation}\label{eq:W}
    W = I + \lambda\, T\,,
\end{equation}
where $I$ represents the BoW contribution (after appropriate normalization) and $T$ is the nearest-neighbor matrix:
\begin{equation}\label{eq:T}
    T_{t,t'} = \delta_{|t-t'|, 1}\,,
\end{equation}
i.e., $T$ is the adjacency matrix of the 1D chain on $L$ sites.

\section{Polar decomposition and dressed embeddings}\label{sec:polar}

\subsection{Absorbing $\sqrt{W}$ into the embedding}

The bilinear form~\eqref{eq:bilinear} can be rewritten as
\begin{equation}
    \tilde{\bm{x}}^\top W\, \tilde{\bm{x}} = \tilde{\bm{x}}^\top \big(\sqrt{W}\big)^\top \sqrt{W}\, \tilde{\bm{x}} = \|\sqrt{W}\,\tilde{\bm{x}}\|^2 = \|\tilde{\bm{z}}'\|^2\,,
\end{equation}
where $\tilde{\bm{z}}' = \sqrt{W}\, \tilde{\bm{x}}$. This defines a new set of \emph{dressed} centered embeddings $\tilde{\bm{z}}'_t$ obtained by mixing the original centered embeddings across context positions.

\subsection{Polar decomposition}

Write $W = U_W\, S_W\, U_W^\top$ (eigendecomposition, since $W$ is symmetric). Then $\sqrt{W} = U_W\, S_W^{1/2}\, U_W^\top$, and the dressed embedding is
\begin{equation}\label{eq:dressed-centered}
    \tilde{\bm{z}}'_t = \sum_{t'} \big(\sqrt{W}\big)_{t,t'}\, \tilde{\bm{x}}_{t'}\,.
\end{equation}
Undoing the centering, this corresponds to a dressed token embedding $\ee_i'$ satisfying
\begin{equation}\label{eq:dressed}
    \boxed{\;\ee_{w_t}' - \bar{\ee}' = \sum_{t'} \big(S_W^{1/2}\, U_W^\top\big)_{t,t'}\, \big(\ee_{w_{t'}} - \bar{\ee}\big)\;}
\end{equation}

\subsection{Perturbative expansion}

For $\lambda \ll 1$:
\begin{equation}\label{eq:sqrt-expand}
    \sqrt{W} = (I + \lambda T)^{1/2} \approx I + \tfrac{1}{2}\lambda\, T + O(\lambda^2)\,.
\end{equation}
So the dressed centered embedding is, to first order:
\begin{equation}\label{eq:dressed-perturbative}
    \tilde{\bm{z}}'_t \approx \tilde{\bm{x}}_t + \frac{\lambda}{2}\big(\tilde{\bm{x}}_{t-1} + \tilde{\bm{x}}_{t+1}\big)\,,
\end{equation}
i.e., each centered embedding receives a correction proportional to the average of its two nearest neighbors in the context window.

\paragraph{Caveat.} The mixing~\eqref{eq:dressed} acts in \textbf{context space} ($L \times L$ matrix on positions), not in token space ($N \times N$ matrix on vocabulary). Since a single token type $i$ can appear at multiple positions, absorbing $\sqrt{W}$ exactly into a position-independent token embedding $\ee_i'$ is not possible in general. We therefore seek the \emph{best approximation} via an optimization.

\section{Optimization: finding the best dressed embedding}\label{sec:optimize}

\subsection{The reconstruction loss}

We want a set of token embeddings $\{\ee_i'\}_{i=1}^N$ that best approximates the context-space dressing. Define the loss
\begin{equation}\label{eq:loss}
    \mathcal{L} = \sum_{t=1}^{L} \big\|\ee_{w_t}' - \big(\sqrt{W}\, \tilde{\bm{x}}\big)_t\big\|^2\,.
\end{equation}

\subsection{Grouping by token identity}

Let $\gamma(i) = \{t : w_t = i\}$ be the set of positions where token $i$ appears, with $|\gamma(i)| \approx (P_i + n_i) \cdot L$. The loss decomposes as
\begin{equation}\label{eq:loss-grouped}
    \mathcal{L} = \sum_{i=1}^{N} \sum_{t \in \gamma(i)} \big\|\ee_i' - \big(\sqrt{W}\, \tilde{\bm{x}}\big)_t\big\|^2\,.
\end{equation}

\subsection{First-order expansion}

Using~\eqref{eq:dressed-perturbative}, the dressed vector at position $t$ is $\tilde{\bm{x}}_t + \frac{\lambda}{2}(\tilde{\bm{x}}_{t-1} + \tilde{\bm{x}}_{t+1})$. Writing $\ee_i' = \ee_i + \delta\ee_i$, the loss at first order in $\lambda$ becomes
\begin{equation}\label{eq:loss-first-order}
    \mathcal{L}(\delta\ee_i) = \sum_{i=1}^{N} \sum_{t \in \gamma(i)} \Big\|\delta\ee_i - \frac{\lambda}{2}\, \ee_{w_{t+1}}\Big\|^2\,,
\end{equation}
where we have used the fact that the dominant nearest-neighbor contribution at each position $t$ is $\frac{\lambda}{2}\, \ee_{w_{t+1}}$ (and similarly $\frac{\lambda}{2}\, \ee_{w_{t-1}}$; by symmetry these contribute equally).

\subsection{Introducing bigram statistics}

For a given token $i$ at position $t$, the token at position $t+1$ is distributed according to the bigram transition probability $P_{ij}^{(b)} = P(w_{t+1} = j \mid w_t = i)$. Replacing the sum over positions $t \in \gamma(i)$ by an average over bigram statistics:
\begin{equation}\label{eq:loss-bigram}
    \mathcal{L}(\delta\ee_i) = \sum_{i=1}^{N} (P_i + n_i) \sum_{j=1}^{N} P_{ij}^{(b)}\, \Big\|\delta\ee_i - \frac{\lambda}{2}\, \tilde{\ee}_j\Big\|^2\,,
\end{equation}
where $\tilde{\ee}_j = \ee_j - \bar{\ee}$ is the centered embedding of token $j$.

\section{Main result: embedding shift as bigram-weighted average}\label{sec:result}

\subsection{Optimization}

Minimizing~\eqref{eq:loss-bigram} with respect to $\delta\ee_i$:
\begin{equation}
    \frac{\partial \mathcal{L}}{\partial(\delta\ee_i)} = 2\sum_j P_{ij}^{(b)} \Big(\delta\ee_i - \frac{\lambda}{2}\, \tilde{\ee}_j\Big) = 0\,.
\end{equation}
Since $\sum_j P_{ij}^{(b)} = 1$, this gives
\begin{equation}\label{eq:shift-raw}
    \delta\ee_i = \frac{\lambda}{2} \sum_j P_{ij}^{(b)}\, \tilde{\ee}_j\,.
\end{equation}

\subsection{The dressed embedding}

\begin{equation}\label{eq:main-result}
    \boxed{\;\ee_i' = \ee_i + \frac{\lambda}{2} \sum_{j=1}^{N} P_{ij}^{(b)}\, \big(\ee_j - \bar{\ee}\big) + O(\lambda^2)\;}
\end{equation}
\textbf{The shift in the embedding of token $i$ is a weighted average of the centered embeddings of all other tokens, with weights given by the bigram transition probabilities $P_{ij}^{(b)} = P(w_{t+1} = j \mid w_t = i)$.}

\subsection{Centering constraint}

The new embedding to first order in $\lambda$ satisfies
\begin{equation}
    \sum_j P_j\, \delta\ee_j = \frac{\lambda}{2} \sum_j P_j \sum_k P_{jk}^{(b)}\, \tilde{\ee}_k\,,
\end{equation}
which ensures the dressed embeddings maintain consistent centering.

\subsection{Self-consistency: bigram probabilities from the model}

The bigram transition probability is determined by the full model~\eqref{eq:perturbed}. Conditioning on $w_t = i$:
\begin{equation}\label{eq:bigram-self}
    P(w_{t+1} = j \mid w_t = i) = \frac{P(w_{t+1} = j,\, w_t = i)}{P(w_t = i)} \approx P_j \cdot e^{-\lambda\, \ee_i \cdot \ee_j}\,,
\end{equation}
where $P_j$ is the unigram frequency and the exponential factor comes from the nearest-neighbor term. At $\lambda = 0$, $P_{ij}^{(b)} = P_j$ (independence), and the shift~\eqref{eq:shift-raw} reduces to $\delta\ee_i = \frac{\lambda}{2}\sum_j P_j\, \tilde{\ee}_j = 0$ (since $\sum_j P_j \tilde{\ee}_j = 0$ by definition of centering). The perturbation is self-consistently non-trivial only when $\lambda \neq 0$.

\section{Open questions}\label{sec:open}

\begin{enumerate}
    \item \textbf{Higher-order expansion and preferred $n$-grams.} The present analysis is first-order in $\lambda$. At higher orders, the dressing involves trigram and higher $n$-gram statistics. Can these corrections be organized systematically? Is there a non-perturbative regime where certain $n$-grams are strongly preferred?

    \item \textbf{Is the tautological interaction program physical?} The perturbation $V = \lambda \sum_t \ee_{w_t} \cdot \ee_{w_{t+1}}$ is constructed from the same embeddings that the Gaussian theory determines---it is ``tautological'' in the sense that the interaction is derived from the unperturbed 2-point function. The self-consistency~\eqref{eq:bigram-self} partially addresses this, but the question remains: does this procedure converge, and does it correspond to a physical (measurable) property of natural language?

    \item \textbf{Computing bigram statistics from the Gaussian theory.} The result~\eqref{eq:main-result} requires the bigram transition probabilities $P_{ij}^{(b)}$, which in principle follow from~\eqref{eq:bigram-self}. But this requires knowing $\ee_i \cdot \ee_j = \Sigma_{ij}$ from the unperturbed theory. The self-consistent loop is: (i)~compute $\Sigma_{ij}$ from the BoW model, (ii)~compute $P_{ij}^{(b)} \approx P_j\, e^{-\lambda \Sigma_{ij}}$, (iii)~compute $\delta\ee_i$ from~\eqref{eq:main-result}. Does this iteration converge?
\end{enumerate}

\end{document}
